\documentclass[11pt,a4paper]{article}
\usepackage[utf8]{inputenc}
\usepackage[T1]{fontenc}
\usepackage[portuguese]{babel}
\usepackage{xcolor}
\usepackage{hyperref}
\usepackage{geometry}
\usepackage{tcolorbox}

\geometry{margin=2.5cm}

\definecolor{primary}{RGB}{41, 128, 185}
\definecolor{secondary}{RGB}{44, 62, 80}
\definecolor{codebg}{RGB}{240, 240, 240}

\hypersetup{colorlinks=true, linkcolor=primary, urlcolor=primary}

\title{\color{secondary}\Huge\textbf{Exercício 5.8: O Panorama (Landscape)}}
\author{\color{primary}DevOps com Kubernetes -- 2026}
\date{}

\begin{document}
\maketitle

\section*{\color{primary}Instruções do Exercício}
\begin{tcolorbox}[colback=blue!5!white,colframe=blue!75!black,title=Exercício 5.8: O Panorama (Landscape)]
Acesse e navegue pelo vasto documento oficial da Cloud Native Computing Foundation: o \textbf{CNCF Cloud Native Landscape} (disponível em PNG ou versão interativa online).

- Em um editor de imagens, circule o logotipo de cada produto ou projeto que você já usou ou tem conhecimento funcional básico (dentro ou fora deste curso).

- Use uma cor diferente para circular todas as \textit{dependências} tecnológicas subjacentes usadas por esses projetos principais.

- Faça uma lista textual anexada detalhando os cenários e contextos de arquitetura em que você precisou deles.
\end{tcolorbox}

\end{document}
